The analog front end is implemented in \code{wheatstone.c}, which owns the bridge excitation and the IADC sampling and hides the pin routing and converter bring-up behind three calls used by \code{app.c}.
The Wheatstone bridge is excited from a single push-pull GPIO output on PC00 and measured differentially by the IADC, with PC01 as the positive input and PC02 as the negative input.
Driving the excitation from a GPIO pin rather than a permanent supply rail allows the bridge to be powered only while a measurement is being prepared, so that it draws no bias current during deep sleep.

The bridge excitation is enabled in \code{wheatstone\_init()} during start-up rather than immediately before the conversion.
By powering the bridge at the start of the cycle, the excitation voltage and the resistive divider have the entire network bring-up window in which to settle, so that the single conversion taken once the network is up samples a stable voltage without requiring an explicit settling delay in the measurement path.
After the conversion the excitation is switched off again, and \code{wheatstone\_prepare\_for\_em4()} guarantees that the supply pin is driven low and the input pins are parked before EM4 entry, which prevents any bridge bias current or floating-pin leakage during sleep.

The IADC is clocked from the 20\,MHz FSRCO, with the analog conversion clock prescaled to approximately 10\,MHz.
The FSRCO is chosen because it starts quickly and does not require a high-frequency crystal to be running, which suits a converter that is enabled only briefly for a single conversion per cycle.
To suppress noise on the small bridge differential signal, the conversion combines 32$\times$ oversampling with an additional 16$\times$ digital averaging, at the cost of a longer conversion time that is acceptable because only one conversion is performed per cycle.
The IADC clocks are gated on only for the duration of the conversion and switched off immediately afterwards, so the converter contributes no current between measurements.

The IADC reports each conversion as an integer rather than a voltage.
It has a 12-bit base resolution, so in differential mode it divides its input range into $2^{12}=4096$ equal steps and returns the step into which the measured voltage falls, where each such integer is referred to as a converter code. 
The span of voltage that this range covers is the differential full-scale input range, which is set by two parameters, namely the reference voltage and the analog gain, where the full-scale range is the reference voltage divided by the analog gain.
The gain is applied to the differential bridge signal by an amplifier stage ahead of the converter, so it scales the analog voltage before it is quantized rather than scaling the resulting integer.
The voltage represented by a single code follows from this, being the total span from the negative to the positive full-scale limit, equal to twice the full-scale range, divided across the 4096 steps.
Two references that are relevant for this purpose are available, the internal 1.2\,V bandgap reference and the 1.8\,V IOVDD supply rail, the latter being the same rail from which the GPIO bridge excitation is driven.
The resulting full-scale ranges for each reference and gain combination are listed in \cref{tab:iadc_gain}, together with the voltage represented by a single converter code and the fraction of each range occupied by the expected bridge swing of approximately $\pm0.2$\,V, which is derived from the bridge dimensioning in \cref{eq:swing_max} and the known resistance change of the sensors from (Section???)

% Pede er nu nået hertil

\begin{table}
\centering
\small
\setlength{\tabcolsep}{4pt}
\caption{IADC differential full-scale range, voltage per converter code at the 12-bit base resolution, and the fraction of range occupied by the expected $\pm0.2$\,V bridge swing, for each analog gain setting and reference voltage.}
\label{tab:iadc_gain}
\begin{tabular}{ccccccc}
\hline
 & \multicolumn{3}{c}{$V_{\mathit{ref}}=1.2$\,V (internal)} & \multicolumn{3}{c}{$V_{\mathit{ref}}=1.8$\,V (IOVDD)} \\
\cmidrule(lr){2-4}\cmidrule(lr){5-7}
Analog gain & Full-scale & V / code & $\pm0.2$\,V uses & Full-scale & V / code & $\pm0.2$\,V uses \\
\hline
$0.5\times$ & $\pm2.4$\,V & \SI{1172}{\micro\volt} & $\sim8\,\%$  & $\pm3.6$\,V & \SI{1758}{\micro\volt} & $\sim6\,\%$ \\
$1\times$   & $\pm1.2$\,V & \SI{586}{\micro\volt}  & $\sim17\,\%$ & $\pm1.8$\,V & \SI{879}{\micro\volt}  & $\sim11\,\%$ \\
$2\times$   & $\pm0.6$\,V & \SI{293}{\micro\volt}  & $\sim33\,\%$ & $\pm0.9$\,V & \SI{439}{\micro\volt}  & $\sim22\,\%$ \\
$4\times$   & $\pm0.3$\,V & \SI{146}{\micro\volt}  & $\sim67\,\%$ & $\pm0.45$\,V & \SI{220}{\micro\volt} & $\sim44\,\%$ \\
\hline
\end{tabular}
\end{table}

The voltage-per-code values in \cref{tab:iadc_gain} are the steps of the 12-bit base resolution.
The 32$\times$ oversampling and 16$\times$ digital averaging do not alter this step directly but lower the noise on the result, so that the smallest reliably resolvable voltage change approaches the quantisation step rather than being masked by noise.

A higher analog gain narrows the full-scale range, so the fixed number of available codes is spread across a smaller voltage span and each code represents a smaller voltage step.
The bridge signal is therefore described by more distinct codes, which means it is quantised more finely and its effective resolution improves.
The penalty is reduced headroom, since a signal that exceeds the narrowed range falls outside the codes the converter can represent and is clipped at the highest or lowest code.
At both reference voltages the highest available gain of 4 keeps the $\pm0.2$\,V swing comfortably inside the range while leaving margin for bridge imbalance, component tolerances, temperature drift, and momentary over-strain, so the gain is set to 4 in either case.

The choice of reference is a trade-off between absolute accuracy and ratiometric stability rather than a question of range alone.
The internal 1.2\,V bandgap reference is more stable in absolute terms and lower in noise, and at a gain of 4 its $\pm0.3$\,V range matches the expected swing most closely, with the signal occupying roughly two-thirds of the range.
The 1.8\,V IOVDD reference instead makes the measurement ratiometric, because the bridge is excited from the same IOVDD rail through the GPIO output, so that drift in the supply scales the bridge output and the reference together and cancels in the result, as discussed for the ratiometric case in \cref{sec:Wheatstone_theory}.
The cost is that the IOVDD rail is noisier than the bandgap reference and its larger full-scale range uses less of the converter span at a given gain, which slightly coarsens the per-sample resolution.
The larger full-scale range would become an advantage if the bridge swing were appreciably greater than $\pm0.2$\,V, since a swing approaching or exceeding $\pm0.3$\,V would clip against the 1.2\,V reference while the wider $\pm0.45$\,V range of the IOVDD reference would still accommodate it.
For the present design, where the swing stays within $\pm0.2$\,V, the firmware uses the internal 1.2\,V reference, prioritising low per-sample noise and absolute accuracy, while the ratiometric IOVDD reference remains a viable alternative where long-term drift over the deployment lifetime is the dominant concern.

The configuration described above is applied in \code{wheatstone\_iadc\_enable()}, the relevant part of which is shown in \cref{lst:iadc_config}.
The default-initialized structures are first populated with the chosen reference, analog gain, oversampling ratio, and digital averaging, after which the IADC is initialized and the differential input pair is routed to PC01 and PC02.

\begin{lstlisting}[language=C, caption={IADC configuration in \code{wheatstone\_iadc\_enable()}, with the clock setup omitted for brevity.}, label={lst:iadc_config}]
IADC_Init_t init = IADC_INIT_DEFAULT;
IADC_AllConfigs_t initAllConfigs = IADC_ALLCONFIGS_DEFAULT;
IADC_InitSingle_t initSingle = IADC_INITSINGLE_DEFAULT;
IADC_SingleInput_t singleInput = IADC_SINGLEINPUT_DEFAULT;

/* ... clock enable, FSRCO select and source prescaler ... */

initAllConfigs.configs[0].reference = iadcCfgReferenceInt1V2;
// 4x analog gain narrows full-scale to ~+-0.3 V for the small bridge signal.
initAllConfigs.configs[0].analogGain = iadcCfgAnalogGain4x;

// 32x oversampling plus digital averaging improves noise for bridge readings.
initAllConfigs.configs[0].osrHighSpeed = iadcCfgOsrHighSpeed32x;
initAllConfigs.configs[0].digAvg = IADC_CFG_DIGAVG_AVG16;

IADC_init(IADC0, &init, &initAllConfigs);

// Differential input: PC1 (pos) - PC2 (neg)
singleInput.posInput = iadcPosInputPortCPin1;
singleInput.negInput = iadcNegInputPortCPin2;
IADC_initSingle(IADC0, &initSingle, &singleInput);
\end{lstlisting}

A measurement is taken by a single call to \code{wheatstone\_read()}, which enables the IADC, issues one single-shot differential conversion across PC01 and PC02, retrieves the result, and then gates the IADC clocks and the bridge excitation back off.
The function returns a signed differential result, where the sign reflects the polarity of the bridge unbalance.